{
Puji syukur ke hadirat Allah SWT atas limpahan rahmat, karunia, serta petunjuk-Nya sehingga penulis dapat menyelesaikan penyusunan skripsi ini dengan baik. Dalam perjalanan menyelesaikan tugas akhir ini, penulis banyak mendapatkan bantuan, dukungan, dan bimbingan dari berbagai pihak. Oleh karena itu, pada kesempatan ini penulis ingin menyampaikan rasa terima kasih kepada:

\begin{enumerate}
	\item Bapak Prof. Ir. Hanung Adi Nugroho, S.T., M.Eng., Ph.D., IPM., SMIEEE., selaku Ketua Departemen Teknik Elektro dan Teknologi Informasi Fakultas Teknik Universitas Gadjah Mada.
	
	\item Bapak Dr. Ahmad Nasikun, S.T., M.Sc., selaku Ketua Program Studi Teknologi Informasi.
	
	\item Bapak Widyawan, S.T., M.Sc., Ph.D., selaku dosen  pembimbing pertama yang banyak membantu dalam memberikan arahan, ide, kritik, serta saran selama proses penyusunan skripsi ini.
	
	\item Bapak Dr. Ir. Guntur Dharma Putra, S.T., M.Sc. selaku dosen pembimbing kedua yang banyak memberikan masukan, saran, serta dukungan selama proses penyusunan skripsi ini.
		
	\item Bapak Prof. Ir. Selo, S.T., M.T., M.Sc., Ph.D., IPU, ASEAN Eng., selaku dosen pembimbing akademik yang telah memberikan bimbingan dan arahan yang sangat berharga selama masa studi penulis di Universitas Gadjah Mada.
	
	\item Seluruh Bapak dan Ibu dosen beserta segenap Tenaga Kependidikan Departemen Teknik Elektro dan Teknologi Informasi (DTETI), yang tidak dapat penulis sebutkan satu per satu, yang telah memberikan ilmu, dedikasi, dan arahan selama masa studi penulis di Universitas Gadjah Mada.

	\item Kedua orang tua penulis, Bapak Fajar Juliadi dan Ibu Muftiatin, yang telah memberikan dukungan, doa, serta kasih sayang yang tiada henti selama penulis menempuh pendidikan hingga saat ini. Tanpa dukungan dan doa dari kedua orang tua, penulis tidak akan dapat menyelesaikan skripsi ini dengan baik.

	\item Saudara penulis, kakak Izza Khorida Bahiya dan Adik Atiqah Alya Fajriah, yang selalu memberikan dukungan, semangat, serta motivasi selama penulis menyelesaikan skripsi ini.

	\item Secara khusus, sahabat penulis, Alm Septian Eka Rahmadi, yang telah menemani dan banyak membantu penulis pada masa studi selama 3 tahun. Semoga segala amal ibadah dan kebaikan beliau diterima di sisi Allah SWT, serta diberikan tempat terbaik disana.

	\item Teman-teman dekat penulis di lingkungan DTETI, yakni Edo, Hanif, Marchel, Brian, Aji, Bayu, Haiqal, Dani, Rafi, Farel, Jemal yang telah memberikan banyak dukungan, motivasi, serta senantiasa menjadi tempat berbagi cerita selama perjalanan masa studi penulis.

	\item Teman-teman kontrakan SMA penulis, yakni Sauki, Haidar, Athoya, Inem, Hanip, Guntur, Ghaffar yang telah menemani penulis selama masa perkuliahan di Yogyakarta.
	
	\item Teman-teman KKN Gemerlap Pinang 2025, yang bersama- sama menciptakan banyak kenangan berharga, memberikan pengalaman hidup yang tak terlupakan, serta menjadi keluarga baru yang saling mendukung satu sama lain.
	
	\item Teman-teman penulis angkatan 2022 DTETI yang telah menjadi teman seperjuangan selama menempuh pendidikan di Universitas Gadjah Mada, serta memberikan banyak pengalaman berharga yang tidak akan terlupakan.


\end{enumerate}

Akhir kata, penulis berharap semoga skripsi ini dapat memberikan manfaat dan kontribusi yang berarti bagi pengembangan ilmu pengetahuan, khususnya dalam bidang teknologi informasi dan kajian keilmuan Islam. Penulis juga berharap penelitian ini dapat menjadi pijakan bagi penelitian-penelitian selanjutnya yang lebih mendalam dan komprehensif.

\begin{flushright}
	\begin{tabular}{c}
		Yogyakarta, 29 Juni 2026 \\
		\vspace{1cm} \\
		Penulis
	\end{tabular}
\end{flushright}